\section{Discussion}
\label{sec:discussion}

\paragraph{Report the operating point, then the curve.} The size sweep is a
case study showing that mAP and deployment behaviour are different measurements.
Three \texttt{yolov8n} seeds and five \texttt{yolov8s} seeds are close on test
mAP ($0.979 \pm 0.003$ vs.\ $0.976 \pm 0.003$, intervals overlapping), but
validation mAP favors s by more than a seed standard deviation and the two
families differ by $6.2\times$ in false alarms at the recall-first point. The
result should motivate a threshold-rule sweep rather than a model ranking claim.
The mechanism is not mysterious: mAP integrates over confidence, and a model can
rank detections well while being badly calibrated near the threshold a user will
pick. The recommendation is concrete and cheap. Report, at minimum: the
validation procedure used to choose a threshold, the threshold itself for every
seed, and precision, recall, and false-positives-per-image at that threshold on
test. Report mAP as well, for comparability, but do not let it be the only row.

\paragraph{Threshold transfer is the hidden variable.} Our five seeds select
thresholds of $0.02$, $0.01$, $0.21$, $0.045$ and $0.125$ for the same
architecture, same data, same schedule. That is not a bug in the search; it is
what happens when a validation set has $181$ images and $7$ negatives. The
consequence is that ``model precision'' pooled over seeds is mostly a statement
about which threshold each seed happened to pick. This is consistent with the
two-stage result: the classifier looks valuable exactly where the threshold
transferred badly, and actively harmful where the selection rule chose a
near-inert configuration, but five seeds with shared components do not identify
that as a mechanism. Any ablation that changes a component downstream of a fragile
threshold will inherit that fragility, and a single pooled mean will hide it.

\paragraph{When is a second stage worth it?} On this evidence, narrowly. The
classifier is defensible if the requirement is recall $\ge 0.96$ with minimum
false alarms, and if the detector's threshold is known to be unreliable. It is
not defensible as a general precision improvement, because a tuned single-stage
threshold beats it at max-$F_2$, and its benefit disappears at the seed whose
threshold transferred well. The cost is real and it is not a rounding error:
$16.5 \pm 1.4$\,ms per image end to end against $6.5$\,ms for the single stage,
with $2.8 \pm 0.6$ crops scored per image and a 95th percentile of $5.3$
(Table~\ref{tab:cost}). For a pipeline that processes every frame of a flight,
that is a $2.5\times$ multiplier bought for a gain that is conditional on the
detector's threshold being unreliable. A deployment decision here is not
``which model is more accurate'' but ``how many false alarms is worth $2.5\times$
the frame budget'', and only the second question has an operational answer.

\paragraph{Leakage is worth measuring, not just fixing.} The contamination-free
comparison is a negative result, and we think negative results of this kind
should be routine. Fixing a split is good practice; quantifying what the fix
changed is what tells the reader whether the original numbers were wrong by
$0.002$ or by $0.2$. Here it is $0.0003$ on the comparable subset. We would not
have known that from the audit alone, and we would have been tempted to describe
the leak as catastrophic.

\paragraph{What the saliency audit bought.} Not much, and we say so rather than
dressing it up. On the $12$ multi-object test images the metric we used is not
even well defined, because a single per-image CAM is reused for every detection;
on the $172$ single-object images the detector's evidence is cleanly
object-localised (median energy share $0.98$, enrichment $\approx90\times$), which
is a statement about localisation, not about correctness. The image-level
baseline's low enrichment is partly a promise of the $20\%$ padding geometry
under which the ground-truth box occupies most of the crop, and the frozen
variant is degenerate on negatives. The audit's one useful output was as a
consistency check on two manually identified label errors. Attribution did not
prove any prediction correct, and the revised interpretation keeps it as a sanity
check rather than a mechanistic explanation.
